\subsection{What Failure of the Reduction Means}

\begin{examplebox}[A singular left block]
For
\[
B=\begin{pmatrix}1&2\\2&4\end{pmatrix},
\]
start with
\[
\left[\begin{array}{cc|cc}
1&2&1&0\\2&4&0&1
\end{array}\right].
\]
After
\[
R_2\leftarrow R_2-2R_1,
\]
the left block contains a zero row:
\[
\left[\begin{array}{cc|cc}
1&2&1&0\\0&0&-2&1
\end{array}\right].
\]
It cannot be reduced to the identity, so $B$ is not invertible.
\end{examplebox}

\begin{interpretationbox}[The method reveals singularity]
A missing pivot is not a defect of Gauss--Jordan elimination. It is evidence that the matrix has no inverse.
\end{interpretationbox}

\begin{summarybox}[Gauss--Jordan procedure]
Form $[A\mid I]$, reduce the left block using elementary row operations, and read the inverse from the right block if the left becomes $I$. If the identity cannot be reached, the matrix is singular.
\end{summarybox}
